\documentclass[UTF8]{ctexart}
\usepackage{amsmath,amssymb}
\usepackage[a4paper,margin=2.5cm]{geometry}
\title{高等代数 1019}
\date{}
\begin{document}
\maketitle

\section*{高等代数 1019}

\subsection*{0. Recall}

代数系统：非空集合 $+$ 运算 $+$ 法则。

\textbf{群}：非空集合 $G$ 上定义运算（加法），满足
\begin{enumerate}
\item 结合律；
\item 有零元；
\item 有负元，
\end{enumerate}
则称 $(G,+)$ 是一个群。进一步，若 $(G,+)$ 还有交换律，则称 $(G,+)$ 是一个 Abel 群。

思考：若 $(G,+)$ 是群，则零元、负元唯一。

\textbf{域}：非空集合 $F$ 上定义两个运算——加法 $+$ 与乘法 $\cdot$，满足
\begin{enumerate}
\item $(F,+)$ 是 Abel 群；
\item $(F\setminus\{0\},\cdot)$ 是 Abel 群；
\item 有分配律，
\end{enumerate}
则称 $(F,+,\cdot)$ 是一个域。

例如 $\mathbb{Q},\mathbb{R},\mathbb{C}$，以及
\[
F_p=\{\bar{0},\cdots,\overline{p-1}\}.
\]

\subsection*{1. 复数的表示方式}

\begin{enumerate}
\item 直角坐标表示：$z=x+iy$；
\item 三角表示：$z=|z|(\cos\theta+i\sin\theta)$；
\item 指数表示：$z=|z|e^{i\theta}$。
\end{enumerate}

若
\[
z_1=x_1+iy_1=|z_1|(\cos\theta_1+i\sin\theta_1)=|z_1|e^{i\theta_1},
\qquad
z_2=x_2+iy_2=|z_2|(\cos\theta_2+i\sin\theta_2)=|z_2|e^{i\theta_2},
\]
则有

加减：
\[
z_1\pm z_2=(x_1\pm x_2)+i(y_1\pm y_2).
\]

乘：
\[
\begin{aligned}
z_1\cdot z_2&=(x_1x_2-y_1y_2)+i(x_1y_2+x_2y_1)\\
&=|z_1||z_2|\bigl(\cos(\theta_1+\theta_2)+i\sin(\theta_1+\theta_2)\bigr)\\
&=|z_1||z_2|e^{i(\theta_1+\theta_2)}.
\end{aligned}
\]

除：
\[
\begin{aligned}
\frac{z_1}{z_2}&=\frac{z_1\bar{z}_2}{z_2\bar{z}_2}
=\frac{x_1x_2+y_1y_2}{x_2^2+y_2^2}+i\frac{x_2y_1-x_1y_2}{x_2^2+y_2^2}\\
&=\frac{|z_1|}{|z_2|}e^{i(\theta_1-\theta_2)}\\
&=\frac{|z_1|}{|z_2|}\bigl(\cos(\theta_1-\theta_2)+i\sin(\theta_1-\theta_2)\bigr).
\end{aligned}
\]

乘方：
\[
z_1^n=|z_1|^ne^{in\theta_1}=|z_1|^n(\cos n\theta_1+i\sin n\theta_1).
\]

开方：
\[
z^{\frac{1}{n}}=\bigl(|z|e^{i\theta}\bigr)^{\frac{1}{n}}
=|z|^{\frac{1}{n}}e^{i\frac{\theta+2k\pi}{n}},
\qquad k=0,1,\cdots,n-1.
\]

下面证明这 $n$ 个数互不相等。设 $0\le p<q\le n-1$，若
\[
\frac{e^{i\frac{\theta+2q\pi}{n}}}{e^{i\frac{\theta+2p\pi}{n}}}
=e^{i\frac{2(q-p)\pi}{n}}=1,
\]
则 $n\mid(q-p)$，与 $0<q-p<n$ 矛盾。于是当 $z\neq0$ 时，$z^{\frac{1}{n}}$ 有 $n$ 个值，即 $x^n=z$ 恰好有 $n$ 个不同的解。

e.g. $z\cdot i$ 是将 $z$ 对应的向量逆时针旋转 $\frac{\pi}{2}$，而 $\frac{z}{i}$ 则是顺时针旋转 $\frac{\pi}{2}$。

平面上有点 $P(x,y)$，将 $\overrightarrow{OP}$ 逆时针旋转 $\theta$ 得到 $\overrightarrow{OP'}$，求 $P'$ 的坐标。设 $P'(x',y')$，二者对应的复数分别为 $z,z'$，则有 $z'=z\cdot e^{i\theta}$，于是
\[
x'+iy'=(x+iy)(\cos\theta+i\sin\theta),
\]
即得
\[
\begin{cases}
x'=\cos\theta\,x-\sin\theta\,y,\\
y'=\sin\theta\,x+\cos\theta\,y.
\end{cases}
\]
写成矩阵，即有
\[
\begin{pmatrix}x'\\y'\end{pmatrix}
=
\begin{pmatrix}\cos\theta & -\sin\theta\\ \sin\theta & \cos\theta\end{pmatrix}
\begin{pmatrix}x\\y\end{pmatrix}.
\]

\subsection*{2. 数域}

定义：复数域的一个非空且元素个数至少为 $2$ 的子集 $K$，如果 $K$ 关于加法、乘法是一个域，则称 $K$ 是一个数域。

等价于：如果 $K$ 对加法、减法、乘法、除法封闭，则 $K$ 是一个数域。

例如：$\mathbb{Q},\mathbb{R},\mathbb{C}$ 是数域，且 $\mathbb{C}$ 是最大的数域；$\mathbb{Z},\mathbb{N}$ 不是数域。

定理：$\mathbb{Q}$ 是最小的数域。

证明：设 $K$ 是一个数域。由于 $K$ 是数域且非空，可取 $0\neq a\in K$。由数域对减法封闭，$0=a-a\in K$；由数域对除法封闭，$1=\frac{a}{a}\in K$。因此 $2=1+1\in K,\cdots,n=(n-1)+1\in K$，故 $\mathbb{N}\subseteq K$。又有 $-n=0-n\in K$，所以 $\mathbb{Z}\subseteq K$。进而 $\forall m,n\in\mathbb{Z}$（$n\neq0$），$\frac{m}{n}\in K$，故 $\mathbb{Q}\subseteq K$。于是 $\mathbb{Q}$ 是最小的数域。证毕。

\end{document}
